\documentclass[10pt,twocolumn]{article}

\usepackage[T1]{fontenc}
\usepackage[utf8]{inputenc}
\usepackage{lmodern}
\usepackage{microtype}
\usepackage[a4paper,margin=0.68in,columnsep=0.25in]{geometry}
\usepackage{amsmath,amssymb}
\usepackage{booktabs}
\usepackage{tabularx}
\usepackage{array}
\usepackage{enumitem}
\usepackage{xcolor}
\usepackage{url}
\usepackage{hyperref}
\usepackage{balance}

\urlstyle{same}
\hypersetup{
  colorlinks=true,
  linkcolor=blue!45!black,
  citecolor=blue!45!black,
  urlcolor=blue!55!black,
  pdftitle={Preference Post-Training Qwen2.5-0.5B},
  pdfauthor={Dheeraj Dhillon}
}

\setlength{\parindent}{1em}
\setlength{\parskip}{0pt}
\setlist{nosep,leftmargin=1.3em}
\newcommand{\model}{\textsc{Qwen2.5-0.5B-Instruct}}
\newcommand{\dataset}{\textsc{HelpSteer3-Preference}}
\newcommand{\trl}{\textsc{TRL}}
\newcolumntype{Y}{>{\raggedright\arraybackslash}X}

\begin{document}

\twocolumn[
\begin{@twocolumnfalse}
\begin{center}
{\LARGE\bfseries Preference Post-Training Qwen2.5-0.5B\par}
\vspace{0.25em}
{\large Supervised Fine-Tuning, Reward Modeling, and Guarded PPO on Diverse, Long-Form HelpSteer3 Data\par}
\vspace{0.45em}
{\normalsize Dheeraj Dhillon\par}
{\small July 2026\par}
\end{center}
\vspace{-0.4em}
\begin{abstract}
This work asks whether a complete reinforcement learning from human feedback
(RLHF) pipeline can remain measurable and inspectable under notebook-scale
compute without reducing preference learning to short or single-domain
outputs. Experiments use \model{} and NVIDIA's heterogeneous \dataset{}.
A sequence-length audit motivates a 4,096-token budget for supervised
fine-tuning (SFT) and reward modeling, reducing observed truncation from about
40\% to about 1\%. One epoch of response-only LoRA SFT supplies a common policy
and reference; a two-epoch Bradley--Terry reward model then reaches 65.62\%
held-out pairwise accuracy.

The final policy is trained with guarded, KL-regularized PPO for 12,032
on-policy rollouts and 188 optimizer updates. The run combines domain-balanced
sampling, quantile-calibrated reward clipping, EOS-aware terminal scoring,
smooth repeated-4-gram shaping, exact resume, and memory-efficient response-only
log-probability computation for 768-token generations. On all 2,017 validation
prompts, PPO remains numerically stable and avoids empty-output collapse, but
the learned reward model prefers it over Base in 817 of 2,017 comparisons
(40.51\%) and over SFT in 974 (48.29\%). PPO nevertheless remains close to the
baselines in length and stopping behavior: its median response is 355 tokens,
its cap-hit rate is 16.31\%, and 16.16\% of responses cross a 25\% repeated
4-gram threshold. The result is therefore presented as an end-to-end systems,
stability, and evaluation study---not as evidence that PPO universally improves
the base instruction model. Complete generations and audit metadata are made
available for direct inspection.
\end{abstract}
\vspace{0.8em}
\end{@twocolumnfalse}
]

\section{Introduction}

RLHF connects next-token language modeling with human judgments of helpfulness,
relevance, correctness, clarity, and safety. A canonical pipeline first adapts
a policy to preferred demonstrations, learns a scalar reward from response
comparisons, and then optimizes the policy while constraining drift from a
supervised reference \cite{stiennon2020,ouyang2022}.

This project deliberately pairs a modest policy and compute budget with a
nontrivial preference distribution. \dataset{} includes general, code, STEM,
and multilingual tasks, long conversational contexts, executable snippets,
and multi-step explanations. The small model makes iteration feasible, but it
also makes sequence truncation, reward-model error, stopping behavior, and
value estimation consequential. The Qwen2.5 family is described in
\cite{qwen25}.

The work began with custom SFT, reward-model, and token-level PPO loops, then
migrated trainer mechanics to Hugging Face \trl{} while retaining repository-
owned preprocessing, model initialization, exact-resume state, reward
calibration, evaluation, and qualitative auditing. The central contribution is
not a state-of-the-art aligned model. It is an inspectable pipeline in which
training behavior and failure analysis are treated as equally important
outputs.

\paragraph{Why SFT an instruction model?}
The starting checkpoint is already instruction-tuned. Continued SFT is used as
domain and preference adaptation, not instruction tuning from scratch. It
moves the policy toward the preferred HelpSteer3 response distribution and
creates a reproducible checkpoint shared by reward modeling, the PPO policy,
and the frozen PPO reference. This controlled initialization makes downstream
changes easier to attribute.

\section{Method}

\subsection{Data and sequence handling}

\dataset{} contains 40,476 pairwise records over general, code, STEM, and
multilingual tasks \cite{wang2025}. Each record contains one conversational
context, two candidate responses, an overall preference in
$\{-3,\ldots,3\}$, and up to three annotator judgments. Tied or invalid pairs
are excluded, leaving 36,264 training and 1,917 reward-model validation pairs.
The generation suite retains all 2,017 validation prompts.

All stages use the Qwen chat template and distinct PAD and EOS tokens.
Completions are preserved through EOS; when an example exceeds 4,096 tokens,
old prompt turns are removed before a final prompt-side token fallback. At a
1,024-token budget, 38.47\% of SFT examples and 40.82\% of reward-model pairs
would truncate. At 4,096 tokens, the observed rates fall to 0.83\% and 1.00\%.
Sequence handling is therefore part of the experiment rather than a cosmetic
configuration choice.

\subsection{Supervised fine-tuning}

For prompt $x$, preferred response $y=(y_1,\ldots,y_T)$, and completion mask
$m_t$, response-only SFT minimizes
\begin{equation}
\mathcal{L}_{\mathrm{SFT}}(\theta)
=-\sum_{t=1}^{T}m_t\log\pi_\theta(y_t\mid x,y_{<t}).
\label{eq:sft}
\end{equation}
Prompt and padding tokens do not contribute to the loss. LoRA
\cite{hu2022} rank-16 adapters are trained in attention and feed-forward
projections while the backbone remains frozen. The merged model supplies the
SFT comparison policy and the common initialization/reference for later
stages.

\subsection{Reward modeling}

The reward model begins from the merged SFT backbone and adds a scalar head.
For preferred $y^+$ and rejected $y^-$ under the same prompt, it minimizes the
Bradley--Terry loss
\begin{equation}
\mathcal{L}_{\mathrm{RM}}(\phi)
=-\log\sigma\!\left(r_\phi(x,y^+)-r_\phi(x,y^-)\right).
\label{eq:rm}
\end{equation}
Only score differences are identifiable; the sign of one isolated score is
not a universal quality boundary. A persisted offset estimated on preferred
demonstrations is subtracted consistently during PPO and evaluation.

\subsection{Guarded KL-regularized PPO}

PPO \cite{schulman2017} receives a terminal reward-model score and token-level
KL shaping against the frozen SFT policy:
\begin{align}
r_t^{\mathrm{KL}}
&=-\beta\left[\log\pi_\theta(y_t\mid s_t)
-\log\pi_{\mathrm{ref}}(y_t\mid s_t)\right],\\
R(x,y)&=\widetilde r_\phi(x,y)+\sum_t r_t^{\mathrm{KL}}.
\end{align}
With generalized advantages $\hat A_t$ and
$\rho_t=\pi_\theta(y_t\mid s_t)/\pi_{\mathrm{old}}(y_t\mid s_t)$, the policy
loss is
\begin{equation}
\mathcal{L}_{\mathrm{PPO}}
=-\mathbb{E}_t\!\left[
\min\{\rho_t\hat A_t,
\operatorname{clip}(\rho_t,1-\epsilon,1+\epsilon)\hat A_t\}
\right].
\label{eq:ppo}
\end{equation}

The final path follows high-impact implementation details from the N+ RLHF/PPO
reproduction \cite{huang2024}: zero policy dropout, behavior probabilities
matched to sampling temperature, a frozen SFT reference, an SFT-initialized
reward model, controlled scalar-head initialization, an RM-initialized critic,
masked advantage whitening, Adam $\epsilon=10^{-5}$, and fixed-length sampling
followed by EOS truncation.

The guarded run adds controls motivated by earlier failures. Reward bounds are
calibrated from 4,096 domain-stratified reward-training pairs. Raw terminal
scores are clipped to the 0.5th and 99.5th percentiles,
$[-2.7507,1.5018]$; a missing EOS receives the same lower-bound score. A smooth
penalty is applied only when repeated token 4-gram mass exceeds the 95th
percentile of preferred responses (0.3084). Reward whitening is disabled so
these calibrated penalties retain their scale, while TRL's masked advantage
whitening remains active. Every batch of 64 contains exactly 16 prompts from
each domain, preventing general prompts from dominating the on-policy data.

\subsection{Memory-efficient long-response execution}

Long-form PPO exposed systems costs that are absent from short demonstrations.
The stock rollout path retained full-vocabulary generation scores, implying a
$64\times768\times151{,}665$ tensor, or about 27.77~GiB in float32, before the
optimizer update. The final path generates token IDs with a rollout-only KV
cache and no generation-score history, then recomputes behavior
log-probabilities in chunks from response-position logits only. It removes
prompt columns masked for every row in each chunk, shares the policy backbone
with the reference by disabling the LoRA adapter, and backpropagates policy and
value losses through separate graphs before one accumulated optimizer step.
These transformations preserve rewards, GAE, clipping, effective batch size,
and sampled trajectories while reducing peak allocated memory to 37.64~GiB on
the completed run.

\begin{table*}[t]
\centering
\small
\caption{Final training profile. ``Token accuracy'' is teacher-forced
next-token accuracy over preferred completion tokens, not task accuracy.}
\label{tab:training}
\begin{tabularx}{0.97\textwidth}{@{}lYYYY@{}}
\toprule
\textbf{Stage} & \textbf{Data and budget} & \textbf{Adaptation} &
\textbf{Optimization} & \textbf{Outcome}\\
\midrule
SFT &
36,264 train; 1,917 val.; 4,096 total tokens &
LoRA $r=16$, $\alpha=32$, dropout 0 &
1 epoch; effective batch 32; LR $5{\times}10^{-6}$ &
Train loss 1.0556; val.\ loss 1.1127; token accuracy 72.02\%\\
Reward model &
36,264 pairs; 1,917 val.\ pairs; 4,096 tokens &
SFT backbone; scalar head; LoRA $r=32$, $\alpha=64$ &
2 epochs; effective batch 64; LR $5{\times}10^{-6}$; centering 0.01 &
Val.\ loss 0.6166; pairwise accuracy 65.62\%; mean margin 0.3578\\
Guarded PPO &
12,032 rollouts; 12,007 unique prompts; 768-token responses &
SFT policy/reference; RM-initialized value model &
188 updates; 64 rollouts/update; 4 PPO epochs; LR $3{\times}10^{-6}$;
$\beta=0.10$ &
Final objective KL 0.7305; update KL 0.000537; EOS 82.81\%; no empty collapse\\
\bottomrule
\end{tabularx}
\end{table*}

\section{Training Results}

\subsection{SFT and reward model}

The one-epoch SFT run is stable. Validation loss changes only from 1.1183 at
step 250 to 1.1127 at step 1,134, while completion-token accuracy remains near
72\%. These metrics do not prove better open-ended responses; they establish a
stable HelpSteer-adapted policy and common reference.

The two-epoch reward model reaches 65.62\% overall pairwise accuracy. Accuracy
increases with preference strength: 61.23\%, 66.33\%, and 71.10\% for strengths
1, 2, and 3. Domain accuracy is uneven: 70.23\% on code, 62.91\% on general,
59.26\% on STEM, and 71.82\% on multilingual examples. It is therefore usable
as a noisy preference signal, but not reliable enough to serve as ground truth,
especially for correctness-sensitive STEM outputs.

\subsection{PPO dynamics}

The guarded run completed its configured budget: 12,032 rollouts and 188
updates. The final response-level objective KL to the frozen SFT reference is
0.7305, or 0.00199 per response token. The old-to-new optimizer update remains
much smaller: approximate KL is 0.000537 and clip fraction is 0.00571 at update
188. Final value loss is 0.1076, down from the large early values produced while
the critic adapted to the terminal reward scale. The final rollout batch has
82.81\% EOS completion, 17.19\% cap hits, mean response length 366.8 tokens,
and no empty responses.

Across the run, later batches are shorter, more likely to terminate, and easier
for the critic to predict than early batches. This is evidence that the guarded
training mechanics worked. It is not evidence that the learned reward objective
became an independent measure of human quality.

\section{Evaluation}

\subsection{Protocol}

Base, SFT, and PPO each generate one response for all 2,017 validation prompts
using the same tokenizer, prompt formatting, 3,072-token prompt limit,
768-token response cap, temperature 0.7, and top-$p=1$. The same epoch-two
reward model scores all prompt-response pairs. Generation is sequential and
resumable at the policy-output level; all pairwise comparisons derive from one
shared response table. Because the PPO training judge is also the evaluator,
learned reward is treated as a diagnostic rather than independent human
preference.

\begin{table*}[t]
\centering
\scriptsize
\caption{Full policy-suite evaluation on 2,017 prompts. Repetition is the
fraction of responses whose repeated word-level 4-gram mass exceeds 25\%. The
behavioral diagnostics are placed before reward/win statistics intentionally.}
\label{tab:evaluation}
\begin{tabular}{@{}lrrrrrrrr@{}}
\toprule
\textbf{Policy} & \textbf{Mean tok.} & \textbf{Median tok.} &
\textbf{Cap hit} & \textbf{Rep-4 $>25\%$} & \textbf{Mean reward} &
\textbf{Three-way wins} & \textbf{Win rate} & \textbf{Empty}\\
\midrule
Base & 373.5 & 333 & 16.21\% & 165 (8.18\%) & 0.2229 & 845 & 41.89\% & 0.00\%\\
SFT & 392.5 & 364 & 18.34\% & 374 (18.54\%) & 0.2037 & 585 & 29.00\% & 0.00\%\\
PPO & 384.1 & 355 & 16.31\% & 326 (16.16\%) & 0.1753 & 563 & 27.91\% & 0.00\%\\
Tie & -- & -- & -- & -- & -- & 24 & 1.19\% & --\\
\bottomrule
\end{tabular}
\end{table*}

\begin{table*}[t]
\centering
\small
\caption{Pairwise learned-reward results. Win rates include ties in the
denominator.}
\label{tab:pairwise}
\begin{tabular}{@{}lrrrrr@{}}
\toprule
\textbf{Comparison} & \textbf{Left wins} & \textbf{Right wins} &
\textbf{Ties} & \textbf{Right win rate} & \textbf{Mean right--left delta}\\
\midrule
Base vs SFT & 1,145 & 849 & 23 & 42.09\% & $-0.0192$\\
Base vs PPO & 1,179 & 817 & 21 & 40.51\% & $-0.0476$\\
SFT vs PPO & 1,018 & 974 & 25 & 48.29\% & $-0.0284$\\
\bottomrule
\end{tabular}
\end{table*}

The learned-reward differences are small in the mean, but the win counts favor
Base. PPO is closest to SFT and exceeds SFT on the STEM subset, yet it does not
beat Base in any of the four domains. This is a mixed alignment outcome rather
than a proxy-judge victory.

\subsection{Stopping, repetition, and qualitative audit}

The guarded PPO response distribution is much closer to Base and SFT than the
earlier, more exploitable PPO checkpoint archived in the technical companion.
In the final suite, 16.16\% of PPO responses cross the 25\% repeated-4-gram
threshold and 5.21\% cross 50\%, compared with 8.18\% and 2.13\% for Base and
18.54\% and 7.09\% for SFT. PPO's 16.31\% cap-hit rate is nearly identical to
Base's 16.21\% and below SFT's 18.34\%.

The automated audit produces review queues rather than semantic verdicts: 9
qualified PPO candidates, 120 strong PPO-loss candidates, 326 repetition-risk
rows, and 155 high-reward repetition mismatches. These sets overlap and are
intended to direct manual inspection. Positive cases include clearer structure,
better coverage of multi-part requests, and more compact answers. Negative
cases include incorrect code, confused factual reasoning, prompt restatement,
and occasional extreme loops that still receive high reward. The response
explorer publishes the full prompt, Base and PPO generations, reward scores,
token counts, EOS/cap state, and repeated 4-gram diagnostics so that these
claims can be checked against concrete text.

\section{Discussion}

The strongest supported result is a systems-and-evaluation result. The project
connects long-context preprocessing, parameter-efficient SFT, pairwise reward
modeling, guarded online PPO, exact-resume execution, full-set generation, and
transparent qualitative audit on a diverse preference dataset. The final PPO
run is numerically healthy and materially less prone to the stopping and
repetition failures that motivated its guardrails. It does not, however,
improve aggregate learned-reward preference over the already instruction-tuned
Base model.

This distinction is central. Stable clipping, low update KL, finite value loss,
and high EOS rate show that the optimizer behaves as intended; they do not
establish that the scalar objective captures every property a human reviewer
cares about. Reward calibration and repetition shaping reduce obvious proxy
exploitation, but they cannot make a 65.62\%-accurate reward model an oracle.
The project therefore treats reward scores, behavioral diagnostics, and direct
response inspection as complementary evidence.

The central limitations are that the same reward model supplies both training
and quantitative evaluation; no blinded multi-rater human study is reported;
one sampled generation per policy does not estimate decoding variance; length
remains entangled with completeness, formatting, and repetition; and a 0.5B
policy/reward model has limited factual and long-context capacity. The final
checkpoint completed a predefined 12,000-episode budget rather than being
selected by an independent human metric. These constraints bound the result,
but they do not invalidate the completed pipeline or the stability analysis.

Future work should prioritize independent evaluation over more PPO updates:
blinded human or strong-judge comparisons on a stratified subset, paired
confidence intervals across repeated generations, reward-model hard negatives
from observed failures, execution-based checks for code, stronger STEM
correctness tests, and multi-metric checkpoint selection. Scaling the policy is
most useful after the evaluator is less exploitable.

\section{Conclusion}

One-epoch continued SFT provides a stable HelpSteer-adapted reference, the
reward model learns a meaningful but imperfect ranking signal, and the guarded
PPO run completes 188 updates over 12,032 long-form rollouts without numerical
or empty-output collapse. Quantile reward bounds, exact-EOS scoring, KL
regularization, repetition shaping, domain balance, exact resume, and
response-only memory optimizations make this the most controlled PPO iteration
in the project. Full evaluation does not show aggregate reward-model
superiority over Base, but it shows a policy whose stopping and repetition
behavior remains usable enough for meaningful side-by-side inspection. The
appropriate conclusion is neither that PPO solved alignment nor that the
training did nothing: the project demonstrates how to build, stabilize, audit,
and honestly report an end-to-end RLHF system under constrained compute.

\balance
\begin{thebibliography}{99}

\bibitem{stiennon2020}
N. Stiennon et al.
\newblock Learning to summarize with human feedback.
\newblock \emph{NeurIPS}, 2020.
\newblock \url{https://arxiv.org/abs/2009.01325}.

\bibitem{ouyang2022}
L. Ouyang et al.
\newblock Training language models to follow instructions with human feedback.
\newblock \emph{NeurIPS}, 2022.
\newblock \url{https://arxiv.org/abs/2203.02155}.

\bibitem{huang2024}
S. Huang, M. Noukhovitch, A. Hosseini, K. Rasul, W. Wang, and L. Tunstall.
\newblock The N+ Implementation Details of RLHF with PPO.
\newblock arXiv:2403.17031, 2024.
\newblock \url{https://arxiv.org/abs/2403.17031}.

\bibitem{wang2025}
Z. Wang et al.
\newblock HelpSteer3-Preference: Open Human-Annotated Preference Data across
Diverse Tasks and Languages.
\newblock arXiv:2505.11475, 2025.
\newblock \url{https://arxiv.org/abs/2505.11475}.

\bibitem{hu2022}
E. J. Hu et al.
\newblock LoRA: Low-Rank Adaptation of Large Language Models.
\newblock \emph{ICLR}, 2022.
\newblock \url{https://arxiv.org/abs/2106.09685}.

\bibitem{schulman2017}
J. Schulman, F. Wolski, P. Dhariwal, A. Radford, and O. Klimov.
\newblock Proximal Policy Optimization Algorithms.
\newblock arXiv:1707.06347, 2017.
\newblock \url{https://arxiv.org/abs/1707.06347}.

\bibitem{qwen25}
Qwen Team.
\newblock Qwen2.5 Technical Report.
\newblock arXiv:2412.15115, 2024.
\newblock \url{https://arxiv.org/abs/2412.15115}.

\bibitem{trl}
L. von Werra et al.
\newblock TRL: Transformer Reinforcement Learning.
\newblock Hugging Face, 2020--2026.
\newblock \url{https://huggingface.co/docs/trl/}.

\end{thebibliography}

\end{document}
